\documentclass[11pt]{article}
\usepackage[margin=1in]{geometry}
\usepackage{booktabs}
\usepackage{hyperref}
\usepackage{enumitem}
\setlist{nosep}

\title{Concept Note: Geometric Transportability for\\Quantum Biosensor Readouts}
\author{Elliot Tower \and Nathan [TBD] \and Ian [TBD, MIT Lincoln Lab]}
\date{July 2026 --- Draft for discussion}

\begin{document}
\maketitle
\thispagestyle{empty}

\section*{Problem}

Quantum biosensors---NV-center nanodiamonds, entangled-photon imagers,
spin-labeled protein detectors---achieve extraordinary sensitivity.
The barrier to clinical adoption is not sensitivity but
\textbf{transportability}: sensor readouts drift across devices,
temperatures, tissue types, and decoherence regimes.  A biomarker
discovered on one instrument in one lab fails to replicate on another.
Every Stage~1 winner of the NIH Quantum Sensing Challenge builds
better sensors; none addresses whether the statistical features
extracted from those sensors are robust enough to transport across
conditions.

\section*{Idea}

We propose a \textbf{deterministic geometric post-processing layer}
that sits downstream of any quantum sensor and extracts
decoherence-robust biomarkers.  The core insight: treat the
high-dimensional sensor output as a point cloud on a manifold, then
select features whose geometric properties are invariant under the
distribution shifts caused by decoherence and calibration drift.

Three existing methods from our prior work compose into this layer:

\begin{enumerate}
  \item \textbf{Transport-stability scoring.}  For each candidate
    biomarker feature, measure whether its value transports across
    held-out decoherence regimes (temperature, T$_2$ relaxation time,
    surface noise).  Features that survive transport are
    decoherence-robust by construction.

  \item \textbf{Ollivier--Ricci curvature summarization.}  Construct a
    nearest-neighbor graph over sensor readouts; compute discrete
    Ricci curvature at each node.  Curvature captures local
    concentration geometry and is stable under moderate perturbation,
    making it a natural decoherence-robust summary statistic.

  \item \textbf{Grassmannian subspace alignment.}  When comparing
    readouts across devices or patients, represent each dataset as a
    subspace on the Grassmannian and measure alignment via principal
    angles.  This gives a calibration-free cross-device similarity
    metric.
\end{enumerate}

The layer is fully deterministic: no learned parameters, no training,
no stochastic optimization.  It takes raw sensor output and returns
transportable features with a quantitative robustness certificate.

\section*{Why this is novel}

\begin{itemize}
  \item \textbf{Unoccupied niche.}  Among the 11 Sensing Stage~1
    winners (and 10 Computing winners), all propose new sensor
    hardware or quantum algorithms.  None addresses the statistical
    validity of the features extracted from quantum sensor output.
  \item \textbf{Sensor-agnostic.}  The geometric layer is agnostic to
    the upstream sensor technology---it applies equally to NV-diamond
    magnetometry, entangled-photon imaging, or spin-labeled protein
    detection.
  \item \textbf{Deterministic = reproducible.}  Pre-registered
    analysis with no learned parameters eliminates a major source of
    replication failure in biomarker discovery.
\end{itemize}

\section*{Validation plan}

We validate on simulated quantum-sensor data with controlled
decoherence, following the same pre-registration discipline we use in
all our work (parameters frozen under a commit SHA before any
experiment runs).

\begin{enumerate}
  \item \textbf{Synthetic NV-diamond readouts.}  Simulate
    fluorescence time-traces under varying T$_2$ times, surface noise
    levels, and temperatures.  Extract candidate biomarker features.
    Measure which features transport across regimes.
  \item \textbf{Published datasets.}  Apply to any publicly available
    quantum-sensor datasets (NV-diamond intracellular thermometry,
    entangled-photon microscopy) if they exist with sufficient
    metadata.
  \item \textbf{Cross-device calibration.}  Demonstrate Grassmannian
    subspace alignment as a zero-parameter calibration transfer
    between simulated devices with different noise profiles.
\end{enumerate}

\section*{Team}

\begin{description}[leftmargin=!,labelwidth=2.5cm]
  \item[Elliot Tower] Geometric methods, transport-stability
    framework, Grassmannian/curvature code, pre-registration protocol,
    write-up.  Independent researcher; prior work on mechanistic
    validity, Ollivier--Ricci curvature in psychiatric comorbidity
    networks, drug-target transportability.
  \item[Nathan] Physics rigor---decoherence models, quantum noise
    characterization, T$_2$/T$_1$ simulation fidelity.
  \item[Ian (MIT-LL)] Hardware realism---real sensor noise profiles,
    NV-diamond or photonic platform grounding, institutional
    credibility in quantum sensing.
\end{description}

\section*{Fit to NIH Qu-BIT}

This targets Sensing Area~2 (``Quantum-enabled Approaches for Early
Detection and Diagnostics'') and Area~3 (``Quantum-enabled Sensing
and Imaging Devices for Diagnostics and Monitoring'').  The geometric
layer is complementary to every Stage~1 winner's hardware---it makes
their sensor output more clinically trustworthy, which is exactly the
translational bottleneck NCATS exists to solve.

\vfill
\noindent\textit{Contact: Elliot Tower,
  \href{mailto:elliot@elliottower.ai}{elliot@elliottower.ai}}

\end{document}
